\documentclass{masterthesis}

\usepackage{hyperref} % links
\usepackage{graphicx}


% macros
\newcommand*\javascript{JavaScript}
\newcommand*\nodejs{Node.js}

\begin{document}

\title{Trace Expressions for Runtime Verification of \nodejs{} Systems}

\author{Luca Franceschini}

\advisor{Davide Ancona}

\examiner{Viviana Mascardi}

\maketitle

\tableofcontents

\chapter{Introduction}

\nodejs{}\footnote{\url{https://nodejs.org/en/}} is a \javascript{}\footnote{Formally specified as the ECMAScript standard \cite{ecmascript}.} runtime environment that runs \javascript{} code outside a browser, as opposed to traditional use of the language for client-side web programming.
In the last few years it made \javascript{} a successful general purpose programming language, with \nodejs{} itself becoming the most popular framework, according to Stack Overflow Developer Survey 2018\footnote{\url{https://insights.stackoverflow.com/survey/2018/\#technology-frameworks-libraries-and-tools}}.

\chapter{Conclusion}

{\em Hello world.}
So far we exploited the HTTP protocol to send observed events to the monitor based on the specification.
This has some performance drawbacks: a new connection is made every time an event is sent, which can happen hundreds of times per second. 
Websocket \cite{websocket} is an interesting alternative as it is expressly designed for continuous exchange of data between two parties and, hence, more efficient than HTTP .
Preliminary tests suggest that this would have a considerable impact on performances. 
  
\bibliographystyle{alphaurl}
\bibliography{bib}

\end{document}

